% Allgemeine Grundlagen für Ordnungen von Profilen; am Bandende eingebunden.
\chapter{Präordnungen und ihre Quotienten}

Eine Präordnung verlangt Reflexivität und Transitivität, aber noch keine
Antisymmetrie. Gegenseitig vergleichbare Elemente werden deshalb zunächst
zusammengefasst. Dieser Abschnitt setzt keine Endlichkeit voraus.

\FormulaDefDeltaK[Präordnung]{%
 \operatorname{PrO}(A,\preceq)\coloneqq
 \bigl(\forall x\in A\;x\preceq x\bigr)\land
 \bigl(\forall x,y,z\in A\;((x\preceq y\land y\preceq z)\to x\preceq z)\bigr)
}{PreorderDef}{%
 \DeltaRow{Mengen}{A\dsep x\dsep y\dsep z}
 \DeltaRow{Relationsoperator}{\preceq}
}

\FormulaDefDeltaK[Symmetrischer Kern einer Präordnung]{%
 x\sim_{\preceq}y\coloneqq(x\preceq y\land y\preceq x)
}{PreorderKernelDef}{%
 \DeltaRow{Mengen}{A\dsep x\dsep y}
 \DeltaRow{Relationsoperatoren}{\preceq\dsep\sim_{\preceq}}
}

\Needspace{9\baselineskip}
\FormulaThmDeltaK[Der symmetrische Kern ist eine Äquivalenzrelation]{%
 \operatorname{PrO}(A,\preceq)\vdash\EqRel{A,\sim_{\preceq}}
}{PreorderKernelEquivalence}{%
 \DeltaRow{Mengen}{A\dsep x\dsep y\dsep z}
 \DeltaRow{Relationsoperatoren}{\preceq\dsep\sim_{\preceq}}
}
\Needspace{6\baselineskip}
\begin{tabproofsplitwide}
\Needspace{5\baselineskip}
\proofpartwide{Reflexivität}
 \proofstepwidestar[1]{\operatorname{PrO}(A,\preceq)}{\rA}
 \proofstepwidestar[2]{x\in A}{\rA}
 \proofstepwidestar[1,2]{x\preceq x}{\FormulaRefAuto{PreorderDef}[def]{1,2}}
 \proofstepwidestar[1,2]{x\sim_{\preceq}x}{\FormulaRefAuto{PreorderKernelDef}[def]{\rAI{3,3}}}
\closeproofpartwide
\Needspace{5\baselineskip}
\proofpartwide{Symmetrie}
 \setcounter{proofstepnr}{4}
 \proofstepwidestar[5]{y\in A}{\rA}
 \proofstepwidestar[6]{x\sim_{\preceq}y}{\rA}
 \proofstepwidestar[6]{x\preceq y\land y\preceq x}{\FormulaRefAuto{PreorderKernelDef}[def]{6}}
 \proofstepwidestar[6]{y\preceq x\land x\preceq y}{\rAI{\rAEb{7},\rAEa{7}}}
 \proofstepwidestar[6]{y\sim_{\preceq}x}{\FormulaRefAuto{PreorderKernelDef}[def]{8}}
\closeproofpartwide
\Needspace{5\baselineskip}
\proofpartwide{Transitivität und Abschluss}
 \setcounter{proofstepnr}{9}
 \proofstepwidestar[10]{z\in A}{\rA}
 \proofstepwidestar[11]{y\sim_{\preceq}z}{\rA}
 \proofstepwidestar[11]{y\preceq z\land z\preceq y}{\FormulaRefAuto{PreorderKernelDef}[def]{11}}
 \proofstepwidestar[1,2,5,6,10,11]{x\preceq z}{\FormulaRefAuto{PreorderDef}[def]{1,2,5,10,\rAI{\rAEa{7},\rAEa{12}}}}
 \proofstepwidestar[1,2,5,6,10,11]{z\preceq x}{\FormulaRefAuto{PreorderDef}[def]{1,10,5,2,\rAI{\rAEb{12},\rAEb{7}}}}
 \proofstepwidestar[1,2,5,6,10,11]{x\sim_{\preceq}z}{\FormulaRefAuto{PreorderKernelDef}[def]{\rAI{13,14}}}
 \proofstepwidestar[1]{\EqRel{A,\sim_{\preceq}}}{\FormulaRefAuto{EqRelDef}[def]{\rUI{\rRI{2,4}},\rUI{\rRI{2,\rUI{\rRI{5,\rRI{6,9}}}}},\rUI{\rRI{2,\rUI{\rRI{5,\rUI{\rRI{10,\rRI{6,\rRI{11,15}}}}}}}}}}
\closeproofpartwide
\end{tabproofsplitwide}

% Der allgemeine Transport wird nach allen Quotientengrundlagen eingefügt.

\Needspace{9\baselineskip}
\FormulaThmDeltaK[Die Präordnung ist mit ihrem Kern verträglich]{%
 \begin{aligned}[t]
 &\operatorname{PrO}(A,\preceq)\dsep x,x',y,y'\in A
 \dsep x\sim_{\preceq}x'\dsep y\sim_{\preceq}y'\vdash{}\\[-2pt]
 &x\preceq y\leftrightarrow x'\preceq y'
 \end{aligned}
}{PreorderKernelCompatibility}{%
 \DeltaRow{Mengen}{A\dsep x\dsep x'\dsep y\dsep y'}
 \DeltaRow{Relationsoperatoren}{\preceq\dsep\sim_{\preceq}}
}
\begin{tabproofwide}
 \proofstepwidestar[1]{\operatorname{PrO}(A,\preceq)}{\rA}
 \proofstepwidestar[2]{x,x',y,y'\in A}{\rA}
 \proofstepwidestar[3]{x\sim_{\preceq}x'}{\rA}
 \proofstepwidestar[4]{y\sim_{\preceq}y'}{\rA}
 \proofstepwidestar[3]{x\preceq x'\land x'\preceq x}{\FormulaRefAuto{PreorderKernelDef}[def]{3}}
 \proofstepwidestar[4]{y\preceq y'\land y'\preceq y}{\FormulaRefAuto{PreorderKernelDef}[def]{4}}
 \proofstepwidestar[7]{x\preceq y}{\rA}
 \proofstepwidestar[1,2,3,7]{x'\preceq y}{\FormulaRefAuto{PreorderDef}[def]{1,2,\rAI{\rAEb{5},7}}}
 \proofstepwidestar[1,2,3,4,7]{x'\preceq y'}{\FormulaRefAuto{PreorderDef}[def]{1,2,\rAI{8,\rAEa{6}}}}
 \proofstepwidestar[1,2,3,4]{x\preceq y\to x'\preceq y'}{\rRI{7,9}}
 \proofstepwidestar[11]{x'\preceq y'}{\rA}
 \proofstepwidestar[1,2,3,11]{x\preceq y'}{\FormulaRefAuto{PreorderDef}[def]{1,2,\rAI{\rAEa{5},11}}}
 \proofstepwidestar[1,2,3,4,11]{x\preceq y}{\FormulaRefAuto{PreorderDef}[def]{1,2,\rAI{12,\rAEb{6}}}}
 \proofstepwidestar[1,2,3,4]{x'\preceq y'\to x\preceq y}{\rRI{11,13}}
 \proofstepwidestar[1,2,3,4]{x\preceq y\leftrightarrow x'\preceq y'}{\rLRI{10,14}}
\end{tabproofwide}

Im Folgenden sind \(Q=\QuotientSet{A}{\sim_{\preceq}}\) und
\([x]=\EqClass{x}{\sim_{\preceq}}{A}\) Abkürzungen.

\FormulaDefDeltaK[Ordnung auf dem Präordnungsquotienten]{%
 p\le_Q q\coloneqq\exists x\in A\,\exists y\in A\,
 (p=[x]\land q=[y]\land x\preceq y)
}{PreorderQuotientOrderDef}{%
 \DeltaPrem{Präordnung}{\operatorname{PrO}(A,\preceq)}
 \DeltaRow{Mengen}{A\dsep Q\dsep p\dsep q\dsep x\dsep y}
 \DeltaRow{Relationsoperatoren}{\preceq\dsep\le_Q}
}

\Needspace{9\baselineskip}
\FormulaThmDeltaK[Vergleich von Klassen durch beliebige Repräsentanten]{%
 \operatorname{PrO}(A,\preceq)\dsep a,b\in A\vdash
 ([a]\le_Q[b]\leftrightarrow a\preceq b)
}{PreorderQuotientComparison}{%
 \DeltaRow{Mengen}{A\dsep Q\dsep a\dsep b\dsep x\dsep y}
 \DeltaRow{Relationsoperatoren}{\preceq\dsep\le_Q}
}
\begin{tabproofwide}
 \proofstepwidestar[1]{\operatorname{PrO}(A,\preceq)}{\rA}
 \proofstepwidestar[2]{a,b\in A}{\rA}
 \proofstepwidestar[1]{\EqRel{A,\sim_{\preceq}}}{\FormulaRefAuto{PreorderKernelEquivalence}[thm]{1}}
 \proofstepwidestar[4]{[a]\le_Q[b]}{\rA}
 \proofstepwidestar[1,4]{\exists x\in A\,\exists y\in A\,([a]=[x]\land[b]=[y]\land x\preceq y)}{\FormulaRefAuto{PreorderQuotientOrderDef}[def]{1,4}}
 \proofstepwidestar[6]{x,y\in A\land[a]=[x]\land[b]=[y]\land x\preceq y}{\rA}
 \proofstepwidestar[1,2,6]{a\sim_{\preceq}x}{\FormulaRefAuto{EqClassRelFromEqual}[thm]{3,2,6}}
 \proofstepwidestar[1,2,6]{b\sim_{\preceq}y}{\FormulaRefAuto{EqClassRelFromEqual}[thm]{3,2,6}}
 \proofstepwidestar[1,2,6]{a\preceq b\leftrightarrow x\preceq y}{\FormulaRefAuto{PreorderKernelCompatibility}[thm]{1,2,6,7,8}}
 \proofstepwidestar[1,2,6]{a\preceq b}{\FormulaRefAuto{P\leftrightarrow Q,Q\vdash P}{9,6}}
 \proofstepwidestar[1,2,4]{a\preceq b}{\rEEStar{5,6,10}}
 \proofstepwidestar[1,2]{[a]\le_Q[b]\to a\preceq b}{\rRI{4,11}}
 \proofstepwidestar[13]{a\preceq b}{\rA}
 \proofstepwidestar[2,13]{\exists x\in A\,\exists y\in A\,([a]=[x]\land[b]=[y]\land x\preceq y)}{\rEIStar{2,\rII,\rII,13}}
 \proofstepwidestar[1,2,13]{[a]\le_Q[b]}{\FormulaRefAuto{PreorderQuotientOrderDef}[def]{1,14}}
 \proofstepwidestar[1,2]{a\preceq b\to[a]\le_Q[b]}{\rRI{13,15}}
 \proofstepwidestar[1,2]{[a]\le_Q[b]\leftrightarrow a\preceq b}{\rLRI{12,16}}
\end{tabproofwide}

\Needspace{9\baselineskip}
\FormulaThmDeltaK[Der Präordnungsquotient ist partiell geordnet]{%
 \operatorname{PrO}(A,\preceq)\vdash\PartOrd{Q,\le_Q}
}{PreorderQuotientPartialOrder}{%
 \DeltaRow{Mengen}{A\dsep Q\dsep p\dsep q\dsep r\dsep x\dsep y\dsep z}
 \DeltaRow{Relationsoperatoren}{\preceq\dsep\le_Q}
}
\Needspace{6\baselineskip}
\begin{tabproofsplitwide}
\Needspace{8\baselineskip}
\proofpartwideindR[Reflexivität]{\operatorname{PrO}(A,\preceq)\dsep p\in Q\vdash p\le_Qp}{\operatorname{PrO}(A,\preceq)\dsep p\in Q\vdash p\le_Qp}[PreorderQuotientReflexive]
 \proofstepwidestar[1]{\operatorname{PrO}(A,\preceq)}{\rA}
 \proofstepwidestar[2]{p\in Q}{\rA}
 \proofstepwidestar[1]{\EqRel{A,\sim_{\preceq}}}{\FormulaRefAuto{PreorderKernelEquivalence}[thm]{1}}
 \proofstepwidestar[1,2]{\exists x\in A\;p=[x]}{\FormulaRefAuto{QuotientRepresentative}[thm]{3,2}}
 \proofstepwidestar[5]{x\in A\land p=[x]}{\rA}
 \proofstepwidestar[1,5]{x\preceq x}{\FormulaRefAuto{PreorderDef}[def]{1,\rAEa{5}}}
 \proofstepwidestar[1,5]{[x]\le_Q[x]}{\FormulaRefAuto{PreorderQuotientComparison}[thm]{1,\rAEa{5},6}}
 \proofstepwidestar[1,5]{p\le_Qp}{\rIE{\FormulaRefAuto{a=b\vdash b=a}[thm]{\rAEb{5}},7}}
 \proofstepwidestar[1,2]{p\le_Qp}{\rEE{4,[5]\vdots8}}
\closeproofpartwideindR
\Needspace{8\baselineskip}
\proofpartwideindR[Transitivität]{\operatorname{PrO}(A,\preceq)\dsep p,q,r\in Q\dsep p\le_Qq\dsep q\le_Qr\vdash p\le_Qr}{\operatorname{PrO}(A,\preceq)\dsep p,q,r\in Q\dsep p\le_Qq\dsep q\le_Qr\vdash p\le_Qr}[PreorderQuotientTransitive]
 \proofstepwidestar[1]{\operatorname{PrO}(A,\preceq)}{\rA}
 \proofstepwidestar[2]{p,q,r\in Q}{\rA}
 \proofstepwidestar[3]{p\le_Qq}{\rA}
 \proofstepwidestar[4]{q\le_Qr}{\rA}
 \proofstepwidestar[1]{\EqRel{A,\sim_{\preceq}}}{\FormulaRefAuto{PreorderKernelEquivalence}[thm]{1}}
 \proofstepwidestar[1,2]{\exists x\in A\;p=[x]}{\FormulaRefAuto{QuotientRepresentative}[thm]{5,2}}
 \proofstepwidestar[1,2]{\exists y\in A\;q=[y]}{\FormulaRefAuto{QuotientRepresentative}[thm]{5,2}}
 \proofstepwidestar[1,2]{\exists z\in A\;r=[z]}{\FormulaRefAuto{QuotientRepresentative}[thm]{5,2}}
 \proofstepwidestar[9]{x\in A\land p=[x]}{\rA}
 \proofstepwidestar[10]{y\in A\land q=[y]}{\rA}
 \proofstepwidestar[11]{z\in A\land r=[z]}{\rA}
 \proofstepwidestar[1,3,9,10]{x\preceq y}{\FormulaRefAuto{PreorderQuotientComparison}[thm]{1,9,10,\rIE{\rAEb{10},\rIE{\rAEb{9},3}}}}
 \proofstepwidestar[1,4,10,11]{y\preceq z}{\FormulaRefAuto{PreorderQuotientComparison}[thm]{1,10,11,\rIE{\rAEb{11},\rIE{\rAEb{10},4}}}}
 \proofstepwidestar[1,3,4,9,10,11]{x\preceq z}{\FormulaRefAuto{PreorderDef}[def]{1,9,10,11,\rAI{12,13}}}
 \proofstepwidestar[1,3,4,9,10,11]{[x]\le_Q[z]}{\FormulaRefAuto{PreorderQuotientComparison}[thm]{1,9,11,14}}
 \proofstepwidestar[1,3,4,9,10,11]{p\le_Qr}{\rIE{\FormulaRefAuto{a=b\vdash b=a}[thm]{\rAEb{9}},\FormulaRefAuto{a=b\vdash b=a}[thm]{\rAEb{11}},15}}
 \proofstepwidestar[1,2,3,4]{p\le_Qr}{\rEE{6,[9]\vdots\rEE{7,[10]\vdots\rEE{8,[11]\vdots16}}}}
\closeproofpartwideindR
\Needspace{8\baselineskip}
\proofpartwideindR[Antisymmetrie]{\operatorname{PrO}(A,\preceq)\dsep p,q\in Q\dsep p\le_Qq\dsep q\le_Qp\vdash p=q}{\operatorname{PrO}(A,\preceq)\dsep p,q\in Q\dsep p\le_Qq\dsep q\le_Qp\vdash p=q}[PreorderQuotientAntisymmetric]
 \proofstepwidestar[1]{\operatorname{PrO}(A,\preceq)}{\rA}
 \proofstepwidestar[2]{p,q\in Q}{\rA}
 \proofstepwidestar[3]{p\le_Qq}{\rA}
 \proofstepwidestar[4]{q\le_Qp}{\rA}
 \proofstepwidestar[1]{\EqRel{A,\sim_{\preceq}}}{\FormulaRefAuto{PreorderKernelEquivalence}[thm]{1}}
 \proofstepwidestar[1,2]{\exists x\in A\;p=[x]}{\FormulaRefAuto{QuotientRepresentative}[thm]{5,2}}
 \proofstepwidestar[1,2]{\exists y\in A\;q=[y]}{\FormulaRefAuto{QuotientRepresentative}[thm]{5,2}}
 \proofstepwidestar[8]{x\in A\land p=[x]}{\rA}
 \proofstepwidestar[9]{y\in A\land q=[y]}{\rA}
 \proofstepwidestar[1,3,8,9]{x\preceq y}{\FormulaRefAuto{PreorderQuotientComparison}[thm]{1,8,9,\rIE{\rAEb{9},\rIE{\rAEb{8},3}}}}
 \proofstepwidestar[1,4,8,9]{y\preceq x}{\FormulaRefAuto{PreorderQuotientComparison}[thm]{1,8,9,\rIE{\rAEb{8},\rIE{\rAEb{9},4}}}}
 \proofstepwidestar[1,3,4,8,9]{x\sim_{\preceq}y}{\FormulaRefAuto{PreorderKernelDef}[def]{\rAI{10,11}}}
 \proofstepwidestar[1,3,4,8,9]{[x]=[y]}{\FormulaRefAuto{EqClassEqualFromRel}[thm]{5,8,9,12}}
 \proofstepwidestar[1,3,4,8,9]{p=q}{\rIE{\FormulaRefAuto{a=b\vdash b=a}[thm]{\rAEb{8}},\FormulaRefAuto{a=b\vdash b=a}[thm]{\rAEb{9}},13}}
 \proofstepwidestar[1,2,3,4]{p=q}{\rEE{6,[8]\vdots\rEE{7,[9]\vdots14}}}
\closeproofpartwideindR
\Needspace{5\baselineskip}
\proofpartwide{Abschluss}
 \proofstepwidestar[1]{\operatorname{PrO}(A,\preceq)}{\rA}
 \proofstepwidestar[2]{p\in Q}{\rA}
 \proofstepwidestar[1,2]{p\le_Qp}{\FormulaRefAuto{PreorderQuotientReflexive}[thm]{1,2}}
 \proofstepwidestar[4]{q\in Q}{\rA}
 \proofstepwidestar[5]{r\in Q}{\rA}
 \proofstepwidestar[6]{p\le_Qq}{\rA}
 \proofstepwidestar[7]{q\le_Qr}{\rA}
 \proofstepwidestar[1,2,4,5,6,7]{p\le_Qr}{\FormulaRefAuto{PreorderQuotientTransitive}[thm]{1,2,4,5,6,7}}
 \proofstepwidestar[9]{q\le_Qp}{\rA}
 \proofstepwidestar[1,2,4,6,9]{p=q}{\FormulaRefAuto{PreorderQuotientAntisymmetric}[thm]{1,2,4,6,9}}
 \proofstepwidestar[1]{\PartOrd{Q,\le_Q}}{\FormulaRefAuto{Partielle Ordnung}[def]{\rUI{\rRI{2,3}},\rUI{\rRI{2,\rUI{\rRI{4,\rUI{\rRI{5,\rRI{6,\rRI{7,8}}}}}}}},\rUI{\rRI{2,\rUI{\rRI{4,\rRI{6,\rRI{9,10}}}}}}}}
\closeproofpartwide
\end{tabproofsplitwide}
